\documentclass[lang=cn,newtx,10pt,scheme=Chinese,citestyle=authoryear,bibstyle=authoryear]{../../Template/elegantbook}

%%%%%%%%%%%%%%%%%%%%%%%%%%%%%%%%%导航区%%%%%%%%%%%%%%%%%%%%%%%%%%%%%%%%%%%%%%%%%

\title{深度学习：基础与概念}
% \subtitle{}

\author{Christopher M. Bishop \& Hugh Bishop}
\institute{中译本}
\date{\today}
\version{ElegantBook-4.5}

\setcounter{tocdepth}{2}

\logo{blank-photo.png}
\cover{cover.png}

\definecolor{customcolor}{RGB}{0,176,183}
\colorlet{coverlinecolor}{customcolor}

% 正文额外使用的宏包和命令
% ================= 隐藏/显示 证明、解答、笔记 的全局开关 =================
% 默认隐藏内容，且不需要点按钮（保留原本的空白占位）
% 如果想恢复原本带按钮的交互效果，只需把下面第一行"% HideEnvContenttrue"注释掉，把第二行"% HideEnvContentfalse"取消注释即可。
\newif\ifHideEnvContent
% \HideEnvContenttrue
\HideEnvContentfalse
% =========================================================================

\makeatletter
\ifHideEnvContent\else
  \bioinfo{全局显示/隐藏证明/解答/笔记}{\AllToggle}% 全局显示/隐藏证明/解答/笔记
\fi
\makeatother

\usepackage{../../Styles/mystyle-elegantbook}% 注意根据文件位置修改模板和宏包的索引。！

% 临时兼容：reference.bib 尚未录入条目，\citep 暂输出为空（待迁移参考文献后移除本宏）
\providecommand{\citep}[1]{}

% 交叉引用设置（这里不引用这个文件自身）
% 注意这里的引用地址不是.sty文件相对于其他被引用的.tex文件地址，而是需要引用的.tex文件相对于其他被引用的.tex地址！

\usepackage{subfiles}

%%%%%%%%%%%%%%%%%%%%%%%%%%%%%%%%%导航区%%%%%%%%%%%%%%%%%%%%%%%%%%%%%%%%%%%%%%%%%

\begin{document}

%%%%%%%%%%%%%%%%%%%%%%%%%%%%正文格式设置%%%%%%%%%%%%%%%%%%%%%%%%%%%%%%%%%%%%%%%%%%

\maketitle

\frontmatter

\tableofcontents% 添加目录

\mainmatter% 将页眉改为预期版本，并重排页码

%%%%%%%%%%%%%%%%%%%%%%%%%%%%%%%%%%正文部分%%%%%%%%%%%%%%%%%%%%%%%%%%%%%%%%%%%%%%%%

\subfile{contents/chapter-01.tex}
\subfile{contents/chapter-02.tex}
\subfile{contents/chapter-03.tex}
\subfile{contents/chapter-04.tex}
\subfile{contents/chapter-05.tex}
\subfile{contents/chapter-06.tex}
\subfile{contents/chapter-07.tex}
\subfile{contents/chapter-08.tex}
\subfile{contents/chapter-09.tex}
\subfile{contents/chapter-10.tex}
\subfile{contents/chapter-11.tex}
\subfile{contents/chapter-12.tex}
\subfile{contents/chapter-13.tex}
\subfile{contents/chapter-14.tex}
\subfile{contents/chapter-15.tex}
\subfile{contents/chapter-16.tex}
\subfile{contents/chapter-17.tex}
\subfile{contents/chapter-18.tex}
\subfile{contents/chapter-19.tex}
\subfile{contents/chapter-20.tex}

%%%%%%%%%%%%%%%%%%%%%%%%%%%%%%%%%%正文部分%%%%%%%%%%%%%%%%%%%%%%%%%%%%%%%%%%%%%%%%

%%%%%%%%%%%%%%%%%%%%%%%%%%%%%%%%%%附录%%%%%%%%%%%%%%%%%%%%%%%%%%%%%%%%%%%%%%%%%%%%

\appendix

\subfile{contents/appendix-A.tex}
\subfile{contents/appendix-B.tex}
\subfile{contents/appendix-C.tex}

%%%%%%%%%%%%%%%%%%%%%%%%%%%%%%%%%%附录%%%%%%%%%%%%%%%%%%%%%%%%%%%%%%%%%%%%%%%%%%%%

%%%%%%%%%%%%%%%%%%%%%%%%%%%%%%%%%%参考文献%%%%%%%%%%%%%%%%%%%%%%%%%%%%%%%%%%%%%%%

\nocite{*}% 显示所有参考文献（原本只显示已引用的参考文献）
\printbibliography[heading=bibintoc, title=\ebibname]

%%%%%%%%%%%%%%%%%%%%%%%%%%%%%%%%%%参考文献%%%%%%%%%%%%%%%%%%%%%%%%%%%%%%%%%%%%%%%

\end{document}